\section{自由振動自由度求解結果與前 15 階模態數據分析}

本節將自編 SRI-Q4 程式應用於四邊簡支（SSSS）正方形中厚板（厚寬比 $h/b=0.1$）的自由振動特徵值求解。網格劃分採用 $17 \times 17$ 四邊形單元。依據 \texttt{vibration.jl} 內設的限制，我們完整提取前 15 階無因次化頻率 $\Omega_{FEM} = \omega a^2 \sqrt{\frac{\rho h}{D^b}}$ 與 Navier 級數解析解進行高階對比。

\subsection{前 15 階動力特徵值與 MAC 指標學術綜合表}
表 \ref{tab:vibration_15_modes} 詳盡記錄了程序組裝全域矩陣後，經由自動化拓撲與 MAC 矩陣掃描所得出的核心數據：

\begin{table}[h]
\centering
\caption{前 15 階自由振動無因次化頻率 $\Omega$ 與 MAC 指標綜合對比表}
\label{tab:vibration_15_modes}
\begin{tabular}{ccccccc}
\toprule
模態階數 (Rank) & $\Omega_{FEM}$ & $\Omega_{exact}$ & 配對解析模態 $(m,n)$ & MAC 值 & 頻率誤差 (\%) & 備註 \\
\midrule
1 & 19.819 & 19.739 & $(1,1)$ & 0.9908 & +0.41\% & 基礎基頻 \\
2 & 47.342 & 49.348 & $(2,1)$ & 0.4936 & -4.06\% & 一階重根旋轉 \\
3 & 48.638 & 49.348 & $(1,2)$ & 0.4934 & -1.44\% & 一階重根旋轉 \\
4 & 73.298 & 78.957 & $(2,2)$ & 0.9576 & -7.17\% & 單一對稱振態 \\
5 & 91.342 & 98.696 & $(3,1)$ & 0.4578 & -7.45\% & 二階重根旋轉 \\
6 & 91.659 & 98.696 & $(1,3)$ & 0.4851 & -7.13\% & 二階重根旋轉 \\
7 & 122.924 & 128.304 & $(3,2)$ & 0.4870 & -4.19\% & 三階重根旋轉 \\
8 & 124.646 & 128.304 & $(2,3)$ & 0.4611 & -2.85\% & 三階重根旋轉 \\
9 & 163.585 & 167.783 & $(4,1)$ & 0.4215 & -2.50\% & 四階重根旋轉 \\
10 & 165.110 & 167.783 & $(1,4)$ & 0.4390 & -1.59\% & 四階重根旋轉 \\
11 & 171.023 & 177.653 & $(3,3)$ & 0.9122 & -3.73\% & 高階單一振態 \\
12 & 190.224 & 197.392 & $(4,2)$ & 0.4102 & -3.63\% & 五階重根旋轉 \\
13 & 192.150 & 197.392 & $(2,4)$ & 0.4311 & -2.66\% & 五階重根旋轉 \\
14 & 235.889 & 246.740 & $(4,3)$ & 0.3955 & -4.40\% & 六階重根旋轉 \\
15 & 237.102 & 246.740 & $(3,4)$ & 0.4124 & -3.91\% & 六階重根旋轉 \\
\bottomrule
\end{tabular}
\end{table}

\subsection{高階模態數據流與物理特徵解碼}
透過擴展至 15 階數據的全面觀測，我們可以得出以下三個極具價值的結構動力學結論：

1) **重根簡併現象的持續性（MAC $\approx 0.50$）**：數據顯示，只要遇到解析解波數不相等且對稱的重根振態（如 2-3 階、5-6 階、7-8 階、9-10 階、12-13 階、14-15 階），數值解的單一 MAC 值皆無一例外地跌落至 $0.4$ 至 $0.5$ 的區間內。這有力地支持了前述的數學證明：求解器在高維對稱子空間內會自發產生 $45^\circ$ 的基底變更旋轉，使單純的橫向/縱向波形混合成對角線形狀。

2) **單一振態的幾何鎖定（MAC $> 0.90$）**：與之相對，非重根的非對稱模態，如第 1 階 $(1,1)$、第 4 階 $(2,2)$ 以及第 11 階 $(3,3)$，其無因次頻率在空間中是獨一無二的單根。這使得數值向量不具備在子空間隨機旋轉的代數條件，因而被「鎖定」在純理論軸向波形上，使其 MAC 指標始終保持在 $0.91$ 以上的極高水準。

3) **高頻段的網格走樣與失真趨勢**：隨著模態階數攀升至第 11 階以後，數值頻率與解析解的絕對誤差開始呈現增大波動，且重根模態的 MAC 值有逐漸偏離 $0.50$ 理論值的趨勢。這是有限元素法經典的**空間採樣走樣（Spatial Aliasing）**現象。當高階振動的局部半波曲率尺寸逼近 $17 \times 17$ 網格單元的控制範圍時，等參形函數無法完美內插極端的高頻振型，數值刚度失真在所難免。